% drafts/sections/methods.tex
\section{Methods}
\label{sec:methods}

\subsection{Unified Manifold Workspace and Coordinate Mapping}
\label{sec:methods:manifold}

\IfFileExists{sections/generated/params.tex}{% Auto-generated by scripts/TexExporter.jl. Do not edit manually.
% Generated UTC: 2026-06-12T18:41:40
\providecommand{\DzMaxMeters}{6.514}
\renewcommand{\DzMaxMeters}{6.514}
\providecommand{\DzMinMeters}{0.009}
\renewcommand{\DzMinMeters}{0.009}
\providecommand{\EnergyFloorThreshold}{1.0e-04}
\renewcommand{\EnergyFloorThreshold}{1.0e-04}
\providecommand{\GridNodes}{33}
\renewcommand{\GridNodes}{33}
\providecommand{\KappaMassBaseline}{1061}
\renewcommand{\KappaMassBaseline}{1061}
\providecommand{\ProfilesCampaignPool}{0}
\renewcommand{\ProfilesCampaignPool}{0}
\providecommand{\ProfilesStableWindow}{0}
\renewcommand{\ProfilesStableWindow}{0}
\providecommand{\ProfilesTotal}{0\xspace}
\renewcommand{\ProfilesTotal}{0\xspace}
\providecommand{\RankCompression}{0}
\renewcommand{\RankCompression}{0}
\providecommand{\RankMedian}{0.00}
\renewcommand{\RankMedian}{0.00}
\providecommand{\RankStd}{0.00}
\renewcommand{\RankStd}{0.00}
\providecommand{\SilhouettePeak}{0.000}
\renewcommand{\SilhouettePeak}{0.000}
\providecommand{\SpectralOrder}{32}
\renewcommand{\SpectralOrder}{32}
\providecommand{\StretchAlpha}{2.5}
\renewcommand{\StretchAlpha}{2.5}
\providecommand{\SvdToleranceFloor}{7.33e-15}
\renewcommand{\SvdToleranceFloor}{7.33e-15}
\providecommand{\TowerMax}{55.00\,\mathrm{m}}
\renewcommand{\TowerMax}{55.00\,\mathrm{m}}
\providecommand{\TowerMaxMeters}{55.00}
\renewcommand{\TowerMaxMeters}{55.00}
\providecommand{\TowerMin}{1.50\,\mathrm{m}}
\renewcommand{\TowerMin}{1.50\,\mathrm{m}}
\providecommand{\TowerMinMeters}{1.50}
\renewcommand{\TowerMinMeters}{1.50}
\providecommand{\TowerObservationLevels}{8}
\renewcommand{\TowerObservationLevels}{8}
\providecommand{\WindowSEarly}{0.000}
\renewcommand{\WindowSEarly}{0.000}
\providecommand{\WindowSIOP}{0.000}
\renewcommand{\WindowSIOP}{0.000}
\providecommand{\WindowSTransitional}{0.000}
\renewcommand{\WindowSTransitional}{0.000}
\providecommand{\alphaStretchBaseline}{2.5}
\renewcommand{\alphaStretchBaseline}{2.5}
}{}

The manuscript methods narrative for the active coordinate transform is generated during the diagnostics pipeline run to keep equations and runtime map parameters synchronized with the executed configuration.

\IfFileExists{sections/generated/transformations_section.tex}{% ==============================================================================
% AUTOMATICALLY GENERATED BY SPECTRALBL REPORT PIPELINE
% GENERATED ON: 2026-06-12T18:41:40.967
% ACTIVE COMPUTATIONAL MAP OBJECT: TanhMap (Riemannian Vertical Manifold)
% CRITICAL: DO NOT MANUALLY AMEND THIS FILE. ALL CHANGES WILL BE OVERWRITTEN.
% ==============================================================================

\subsection{Analytical Coordinate Mapping and Riemannian Metric Formulations}
To project non-uniformly spaced atmospheric tower instrumentation from CASES-99 into a stable workspace for high-order pseudospectral methods, the pipeline instantiates a curvilinear coordinate chart modeled as a one-dimensional Riemannian manifold.

The physical coordinate $z \in [z_{\min}, z_{\max}]$ is linked to the computational coordinate $\xi \in [-1, 1]$ through smooth invertible maps $z = \mathcal{G}(\xi)$ and $\xi = \mathcal{F}(z)$.

\subsection{Symmetric Hyperbolic Tangent Transformation (\texttt{TanhMap})}
For centered shear-layer transitions, a symmetric hyperbolic tangent map with $\alpha = 2.5$ is applied.
The geometric midpoint is $z_c = 28.25\,\mathrm{m}$ and the span is $L = 53.5\,\mathrm{m}$.

\begin{equation}
    \xi = \mathcal{F}(z) = \frac{1}{\alpha} \operatorname{atanh}\left( \frac{2(z - z_c)\tanh(\alpha)}{L} \right)
\end{equation}
\begin{equation}
    z = \mathcal{G}(\xi) = z_c + \frac{L}{2} \frac{\tanh(\alpha \xi)}{\tanh(\alpha)}
\end{equation}
\begin{equation}
    \frac{dz}{d\xi} = \frac{L}{2} \left( \frac{\alpha}{\tanh(\alpha)} \right) \operatorname{sech}^2(\alpha \xi)
\end{equation}
\begin{equation}
    \frac{d^2z}{d\xi^2} = -L \left( \frac{\alpha^2}{\tanh(\alpha)} \right) \operatorname{sech}^2(\alpha \xi)\tanh(\alpha \xi)
\end{equation}


\subsubsection{Riemannian Metric Tensor and Volume Form}
The covariant and contravariant metric components are
\begin{equation}
    g_{\xi\xi} = J(\xi)^2 = \left(\frac{dz}{d\xi}\right)^2, \qquad
    g^{\xi\xi} = J(\xi)^{-2} = \left(\frac{dz}{d\xi}\right)^{-2}
\end{equation}
where $J(\xi)$ is the local Jacobian metric coefficient. The line element is $ds^2 = g_{\xi\xi} d\xi^2$.

Spatial variations in grid density are governed by the metric-gradient term
\begin{equation}
    J_{\xi}(\xi) = \frac{dJ}{d\xi} = \frac{d^2z}{d\xi^2}
\end{equation}
Because a one-dimensional manifold is intrinsically flat, $J_{\xi}(\xi)$ is a coordinate-stretching gradient, not intrinsic curvature. The induced volume form is
\begin{equation}
    dV = \sqrt{|g|} d\xi = J(\xi) d\xi
\end{equation}

\subsubsection{Geometric Inner Product and Hilbert Space Workspace}
The flow fields are represented in weighted $L^2_g([-1,1])$. For square-integrable fields $f$ and $g$,
\begin{equation}
    \langle f, g \rangle_g = \int_{-1}^{1} f(\xi) g(\xi) dV = \int_{-1}^{1} f(\xi) g(\xi) J(\xi) d\xi
\end{equation}
with induced norm $\|f\|_g^2 = \langle f, f \rangle_g$.

\subsubsection{Pseudospectral Operator Matrix Discretization}
Differentiation operators are assembled on the canonical Chebyshev computational domain using collocation nodes $\xi_j = \cos(\pi j / N)$ and then mapped to physical space:
\begin{equation}
    \frac{\partial \phi}{\partial z} = J^{-1} \frac{\partial \phi}{\partial \xi}
\end{equation}
\begin{equation}
    \frac{\partial^2 \phi}{\partial z^2} = J^{-2} \frac{\partial^2 \phi}{\partial \xi^2} - J^{-3} J_{\xi} \frac{\partial \phi}{\partial \xi}
\end{equation}

With $\mathbf{D}_{\xi} \in \mathbb{R}^{(N+1)\times(N+1)}$ and diagonal projections
\begin{equation}
    \mathbf{J}^{-1} = \operatorname{diag}\left(J^{-1}(\xi_j)\right), \quad
    \mathbf{J}^{-2} = \operatorname{diag}\left(J^{-2}(\xi_j)\right), \quad
    \mathbf{J}^{-3} = \operatorname{diag}\left(J^{-3}(\xi_j)\right)
\end{equation}
the discrete physical operators are
\begin{equation}
    \mathbf{D}_z = \mathbf{J}^{-1} \mathbf{D}_{\xi}
\end{equation}
\begin{equation}
    \mathbf{D}_{zz} = \mathbf{J}^{-2} \mathbf{D}_{\xi}^2 - \mathbf{J}^{-3} \mathbf{J}_{\xi} \mathbf{D}_{\xi}
\end{equation}
where $\mathbf{J}_{\xi} = \operatorname{diag}\left(J_{\xi}(\xi_j)\right)$.

\subsubsection{Geometric Quadrature Integration Consistency}
Continuous integrals transform as
\begin{equation}
    \int_{z_{\min}}^{z_{\max}} \phi(z) dz = \int_{-1}^{1} \phi\left(\mathcal{G}(\xi)\right) J(\xi) d\xi
\end{equation}
For Chebyshev-Gauss-Lobatto computational weights $w_j^{(\xi)}$, the physical weights are
\begin{equation}
    w_j^{(z)} = J(\xi_j) w_j^{(\xi)}
\end{equation}
This enforces dual consistency between metric-weighted differentiation and quadrature summations.
}{}

The active map metadata table is also generated from runtime state and included below for reproducibility.

\IfFileExists{sections/generated/table_transform_config.tex}{% Auto-generated by scripts/TexExporter.jl. Do not edit manually.
\begin{table}[htbp]
  \centering
  \caption{Coordinate transform configuration used for this run.}
  \label{tab:transform-config}
  \begin{tabular}{l l}
    \toprule
    Parameter & Value \\
    \midrule
    map\_type & TanhMap \\
    zmin & 1.5 \\
    zmax & 55.0 \\
    alpha & 2.5 \\
    \bottomrule
  \end{tabular}
\end{table}
}{}

\subsection{Scale Separation and Reduced-Order Representation}
\label{sec:methods:scale_separation}

To avoid conflating empirical model behavior with claims about the full infinite-dimensional atmospheric operator, we explicitly introduce a reduced-order fast/slow representation motivated by observed SBL time scales. We define a fast turbulence-adjustment scale
\begin{equation}
\tau_t \sim \frac{\ell}{u_*},
\end{equation}
where $\ell$ is a local mixing-length scale and $u_*$ is friction velocity, and a slow mesoscale forcing scale $\tau_m$ associated with radiative cooling, geostrophic modulation, and low-level-jet evolution. In strongly stable windows these scales are typically separated, with $\tau_t \sim 10^{1}$--$10^{2}\,\mathrm{s}$ and $\tau_m \sim 10^{4}$--$10^{5}\,\mathrm{s}$.

We therefore use the dimensionless ordering parameter
\begin{equation}
\epsilon = \frac{\tau_t}{\tau_m}, \qquad \epsilon \ll 1,
\end{equation}
as a modeling guide for reduced-state interpretation. This manuscript does not claim that the full Boussinesq/Navier--Stokes system has been rigorously transformed into a singularly perturbed normal form; rather, we use this scale separation to justify a reduced-order diagnostic representation for the observed trajectory geometry.


\subsection{Thin SVD-Truncated Matrix Projection Architecture}
\label{sec:methods:svd}

Reconstructing high-fidelity continuous profiles from a highly constrained, sparse array of vertical observations ($M = \TowerObservationLevels$ physical tower levels) onto an $N=\SpectralOrder$ order spectral expansion requires resolving an inherently rank-deficient linear inverse problem. Evaluating the basis functions at the fixed instrument elevations yields the rectangular evaluation design matrix $\mathbf{H} \in \mathbb{R}^{M \times (N+1)}$:
\begin{equation}
H_{ij} = T_{j-1}(\xi(z_i))
\label{eq:h_matrix_def}
\end{equation}

To minimize computational overhead while guaranteeing protection against numerical over-fitting, we reject classic normal equations or unconstrained Tikhonov diagonal dampening. Instead, we compute an economy-size, thin Singular Value Decomposition (SVD) of the rectangular operator:
\begin{equation}
\mathbf{H} = \mathbf{U}_r \mathbf{S}_r \mathbf{V}_r^T
\label{eq:thin_svd}
\end{equation}
where $\mathbf{U}_r \in \mathbb{R}^{M \times M}$ and $\mathbf{V}_r \in \mathbb{R}^{(N+1) \times M}$ retain strictly orthogonal column components, and $\mathbf{S}_r = \operatorname{diag}(s_1, s_2, \ldots, s_M)$ holds the ordered singular spectrum. This thin variant restricts matrix multiplications exclusively to the data-constrained subspace, eliminating the numerical costs associated with tracking unresolvable high-frequency spectral components.

To isolate valid physical modes from background precision noise without introducing hardcoded scaling artifacts during intense thermal modifications, we enforce a dynamic truncation tolerance criterion $\tau$:
\begin{equation}
\tau = \max(M, N+1) \cdot s_1 \cdot \epsilon_{\mathrm{mach}}
\label{eq:dynamic_tolerance}
\end{equation}
where $\epsilon_{\mathrm{mach}} \approx 2.22 \times 10^{-16}$. The effective operational rank is bounded strictly by $r_{\mathrm{eff}} = \#\{s_i > \tau\}$. The optimized, noise-filtered spectral coefficient vector $\mathbf{c} \in \mathbb{R}^{N+1}$ is computed directly via projection onto the right-singular basis:
\begin{equation}
\mathbf{c} = \sum_{i=1}^{r_{\mathrm{eff}}} \frac{\mathbf{u}_i^T \mathbf{a}_{\mathrm{obs}}}{s_i} \mathbf{v}_i
\label{eq:svd_reconstruction}
\end{equation}
Projections are guaranteed to lie entirely within the data-validated coordinate workspace, automatically forcing unconstrained higher-order spectral modes cleanly to zero.

This rank-truncated representation is paired with event-oriented interpretation: we track localized structural deformations in reconstructed profiles as state transitions, following non-stationary event-detection logic \citep{Kang2014a} and the turbulence-events viewpoint of \citet{Narasimha1995}, rather than assuming that continuous gradient metrics alone provide sufficient state identification.

\subsubsection{Representation Invariance vs. Estimator Uncertainty}
The manifold state definition is coordinate-agnostic with respect to sensor hardware topology: for any admissible tower geometry, the reconstruction is posed in the same basis and metric space, and the diagnostics ($D_{\mathrm{eff}}$, $F_W$, $\chi_N$) retain the same mathematical interpretation. However, \emph{estimation uncertainty is not geometry-agnostic}. Let measurement perturbations satisfy $\mathbf{a}_{\mathrm{obs}} \mapsto \mathbf{a}_{\mathrm{obs}} + \delta \mathbf{a}$. In the truncated subspace, coefficient perturbations scale as
\begin{equation}
\delta \mathbf{c} \approx \mathbf{V}_r \mathbf{S}_r^{-1} \mathbf{U}_r^T \, \delta \mathbf{a}
\label{eq:svd_uncertainty_map}
\end{equation}
so that $\|\delta \mathbf{c}\|_2 \le s_{r_{\mathrm{eff}}}^{-1}\|\delta \mathbf{a}\|_2$. When these mapped perturbations are projected onto physical spatial gradients via the Riemannian differentiation operator $\mathbf{D}_z = \mathbf{J}^{-1}\mathbf{D}_{\xi}$, the spatial distribution of uncertainty is non-uniformly modulated across the column:
\begin{equation}
\|\delta \mathbf{D}_z \mathbf{\phi}\|_2 \le \|\mathbf{J}^{-1}\|_2 \, \|\mathbf{D}_{\xi} \mathbf{V}_r \mathbf{S}_r^{-1} \mathbf{U}_r^T\|_2 \, \|\delta \mathbf{a}\|_2
\end{equation}
Sparse vertical gaps or clustered level placement steepen singular-value decay (smaller $s_{r_{\mathrm{eff}}}$), amplifying uncertainty in the recovered gradients---particularly near regions where the local Jacobian scale factor $J(\xi)$ compresses physical node distributions.

For homoscedastic observational noise with covariance $\mathbf{\Sigma}_a = \sigma_a^2\mathbf{I}$, first-order propagation yields
\begin{equation}
\mathbf{\Sigma}_c \approx \sigma_a^2\, \mathbf{V}_r\, \mathbf{S}_r^{-2}\, \mathbf{V}_r^T
\label{eq:svd_covariance}
\end{equation}
which makes explicit that uncertainty concentration is controlled by the inverse-squared singular spectrum. Accordingly, cross-platform comparisons should treat manifold coordinates as transferable state variables while reporting geometry-conditioned confidence via $(r_{\mathrm{eff}}, s_{r_{\mathrm{eff}}}, \kappa)$ and associated uncertainty envelopes.


\subsection{Smooth Scale Partitioning and Energy Non-Conservation Caveat}
\label{sec:methods:partition}

To separate multi-scale structures without inducing non-local Gibbs ringing or artificial step discontinuities, we deploy a partition-of-unity windowing framework. The global flow field is decomposed into mean synoptic ($\psi_M$), internal gravity wave ($\psi_W$), and micro-turbulent residual ($\psi_T$) scales using coupled hyperbolic tangent transition channels evaluated over physical height $z$:
\begin{equation}
\begin{aligned}
  \psi_T(z) &= \frac{1}{2}\left[1 - \tanh\left(\frac{z - z_{\mathrm{surf}}}{\Delta_z}\right)\right], \\
  \psi_M(z) &= \frac{1}{2}\left[1 + \tanh\left(\frac{z - z_{\mathrm{top}}}{\Delta_z}\right)\right], \\
  \psi_W(z) &= 1 - \psi_M(z) - \psi_T(z)
\end{aligned}
\label{eq:partition_scales_def}
\end{equation}
where the physical transition anchors are structurally mapped to the observation window bounds. By design, the partition operators sum identically to unity across the physical column:
\begin{equation}
\psi_M(z) + \psi_W(z) + \psi_T(z) = 1 \quad \forall z \in [z_{\min}, z_{\max}]
\label{eq:partition_unity}
\end{equation}

Crucially, because these smooth coordinate windows are not projections onto orthogonal eigenvectors of the metric mass matrix $\mathbf{M}$, \emph{global physical energy content is not strictly additive across the partitioned sub-domains}. In discrete form, the windowed physical profiles are $\mathbf{a}^{(j)} = \operatorname{diag}(\psi_j(z_i))\mathbf{a}$ for $j\in\{M,W,T\}$, and their corresponding spectral coefficients are recovered by the same inverse map used for the full field. The cross-product metrics therefore do not vanish identically:
\begin{equation}
\langle \mathbf{c}, \mathbf{M} \mathbf{c} \rangle = \sum_{j \in \{M,W,T\}} \langle \mathbf{c}^{(j)}, \mathbf{M} \mathbf{c}^{(j)} \rangle + \sum_{j \neq k} \langle \mathbf{c}^{(j)}, \mathbf{M} \mathbf{c}^{(k)} \rangle
\label{eq:energy_overlap}
\end{equation}
The off-diagonal cross-terms $\langle \mathbf{c}^{(j)}, \mathbf{M} \mathbf{c}^{(k)} \rangle \neq 0$ physically represent overlapping scale support within the boundary layer transitions. While this cross-term behavior accurately mirrors the coupled wave-turbulence interaction concepts emphasized by \citet{Sun2015}, and is consistent with observed solitary-event and microfront dynamics in very stable flows \citep{Mahrt2010}, total energy balance estimates derived via this windowing scheme must explicitly track these non-orthogonal interaction components rather than assuming direct linear superposition.

These non-vanishing terms are not residual numerical noise: they are the mathematically explicit channels through which wave-support and turbulent-support subspaces exchange energy during regime transition intervals. Operationally, neglecting the off-diagonal block in Eq.~(\ref{eq:energy_overlap}) biases partition diagnostics toward false additivity and suppresses the coupled interaction regions that dominate intermittent shear states.
In implementation, spikes in the diagnosed Interaction Residual are therefore interpreted as physically meaningful submesoscale coupling events (wave--turbulence exchange episodes) rather than as stochastic solver noise \citep{Sun2015, Mahrt2010}.


\subsection{Information-Theoretic Diagnostic Features}
\label{sec:methods:diagnostics}

\subsubsection{Effective Modal Dimension (\texorpdfstring{$D_{\mathrm{eff}}$}{Deff}) with Minimum Energy Filtering}

We compute the effective modal diagnostic $D_{\mathrm{eff}}$ as the exponential of the spectral Shannon entropy, converting a raw information metric into an interpretable proxy for the number of actively participating, equiprobable spectral modes:
\begin{equation}
D_{\mathrm{eff}} = \exp\left( -\sum_{n=0}^{N} p_n \log p_n \right)
\label{eq:d_eff_def}
\end{equation}
where $p_n = c_n^2 / \sum_m c_m^2$ is the normalized energy in mode $n$, applying the standard information-theoretic convention that $p_n \log p_n \to 0$ when $p_n \to 0$.

Because $D_{\mathrm{eff}}$ operates as a scale-invariant shape diagnostic, a quiescent profile under weak background conditions could mathematically register an artificially inflated modal dimension due to numerical precision variations. To prevent low-amplitude instrumental noise from triggering false high-dimensional turbulence classifications, we introduce a regularized minimum energy masking operator before computing the spectral probabilities:
\begin{equation}
c_n^2 \leftarrow \begin{cases}
c_n^2, & \text{if } \sum_{m=0}^{N} c_m^2 \ge \mathcal{E}_{\text{floor}} \\
0, & \text{otherwise}
\end{cases}
\label{eq:energy_floor}
\end{equation}
where the threshold floor is defined dynamically based on historical instrument variance limits, fixed at $\mathcal{E}_{\text{floor}} = 10^{-4} \cdot \max_{t} \left(\|\mathbf{a}_{\mathrm{obs}}\|^2\right)$. This restriction forces highly quiescent, laminar configurations cleanly down to $D_{\mathrm{eff}} \to 1.0$.

In the reduced-order interpretation, low $D_{\mathrm{eff}}$ is consistent with strong contraction toward a wave-dominated neighborhood, while transient inflation of $D_{\mathrm{eff}}$ is interpreted as a proxy for local relaxation of restoring rates near fold-proximal trajectories. We treat this as an empirical indicator of distance from normal-hyperbolicity-like behavior in the reduced chart, not as a direct measurement of manifold dimension.

\subsubsection{Wave Energy Fraction (\texorpdfstring{$F_W$}{FW}) and Spectral Curvature (\texorpdfstring{$\chi_N$}{ChiN})}

The wave energy fraction isolates metric-weighted energy in the wave band relative to total energy:
\begin{equation}
F_W = \frac{\langle \mathbf{c}^{(W)}, \mathbf{M} \mathbf{c}^{(W)} \rangle}{\langle \mathbf{c}, \mathbf{M} \mathbf{c} \rangle}
\label{eq:f_w}
\end{equation}

Spectral curvature ($\chi_N$) measures the concentration of gradient energy in high-frequency modes, providing a normalized fourth-moment proxy for structural profile sharpness:
\begin{equation}
\chi_N = N^{-2} \left( \frac{\sum_{n=1}^{N} n^4 c_n^2}{\sum_{n=1}^{N} n^2 c_n^2} \right)
\label{eq:chi_n}
\end{equation}

\subsubsection{Dual-Sided Spectral Coupling of \texorpdfstring{$D_{\mathrm{eff}}$}{Deff} and \texorpdfstring{$\chi_N$}{ChiN}}
The expected negative coupling between $D_{\mathrm{eff}}$ and $\chi_N$ follows directly from the structure of Eqs.~(\ref{eq:d_eff_def}) and (\ref{eq:chi_n}): the two diagnostics encode complementary limits of the same modal energy distribution. $D_{\mathrm{eff}}$ increases as modal probabilities flatten across many active modes, while $\chi_N$ amplifies high-order tail occupancy through the $n^4$ numerator relative to the bulk gradient-energy scale in the $n^2$ denominator.

In quiescent or weakly forced intervals, modal support collapses toward low-order coefficients, driving $D_{\mathrm{eff}} \to 1$. In that limit, the denominator $\sum n^2 c_n^2$ can remain small, so even weak residual high-order content is strongly magnified by the $n^4$ weighting and $\chi_N$ can rise. During fully developed, broadband mixing, energy spreads across many modes, increasing $D_{\mathrm{eff}}$ while simultaneously normalizing the curvature ratio through a larger gradient-energy denominator, which tends to reduce or stabilize $\chi_N$.

This dual-sided response is advantageous for state classification because the feature pair is physically coupled but not redundant. In the two-dimensional diagnostic space $(D_{\mathrm{eff}}, \chi_N)$, laminar or wave-dominated states occupy a high-$\chi_N$/low-$D_{\mathrm{eff}}$ sector, intermittent transitions align along an oblique intermediate manifold, and continuously turbulent states cluster toward low-$\chi_N$/high-$D_{\mathrm{eff}}$. Consequently, mixture-model clustering is expected to recover a diagonally structured covariance aligned with the laminar-to-turbulent transition pathway rather than axis-aligned degeneracy.

For stability diagnostics, $\mathrm{Ri}_g$ remains the primary manifold-coordinate reference because it uses resolved near-surface gradients from the reconstructed profiles. In parallel, we track the bulk Richardson number $\mathrm{Ri}_b$ as a supplemental finite-difference diagnostic over top-minus-bottom layer means to provide a geometry-aware bridge to tower-layer formulations.